\chapter{Simulation and Numerical Tools for Performance Evaluation}
\label{sec:appendix_sim_tools}

Performance evaluation of medium access control (MAC) protocols and wireless network architectures relies on a tripartite methodology comprising analytical modeling, computer simulation, and experimental prototyping \cite{kumar2004communication, bertsekas1992data}. While analytical models (such as Discrete-Time Markov Chains, renewal theory, and fluid approximations developed throughout this monograph) provide fundamental physical insights and closed-form scaling laws, their mathematical tractability often necessitates simplifying assumptions—such as independent packet losses, Poisson arrivals, or homogeneous spatial configurations. Conversely, experimental testbeds provide absolute physical fidelity but are constrained by hardware scalability, high equipment costs, and environmental non-repeatability.

Computer simulation bridges this gap by enabling high-fidelity, repeatable, and scalable evaluation of complex, heterogeneous networking scenarios under realistic radio propagation, non-stationary traffic, and multi-tier protocol stacks \cite{fall2003ns, henderson2008network}. This appendix provides an extensive reference guide to the premier simulation and numerical computation tools utilized in modern wireless network research:
\begin{enumerate}
    \item \textbf{MATLAB and Simulink:} Matrix-oriented numerical analysis, transient differential equation solvers, Monte Carlo protocol simulators, and PHY-MAC toolboxes.
    \item \textbf{Network Simulator 2 (NS-2):} The classical dual-language (C++/OTcl) discrete-event simulation engine for packet-level wireless protocol stacks.
    \item \textbf{NetSim (Tetcos):} Commercial discrete-event simulation and emulation platform with protocol engines across Wi-Fi, 5G NR, TSCH, and IoT.
    \item \textbf{SimPy:} Modern Python-based process-oriented discrete-event simulation framework, highlighting the MAC evaluation methodology developed by Hemanth, Venkatesh, and Iyer \cite{hemanth2010simpy}.
    \item \textbf{Network Simulator 3 (NS-3):} Architectural paradigms of modular C++ discrete-event network simulation.
\end{enumerate}

---

\section{Fundamentals of Discrete-Event Network Simulation}
\label{sec:des_fundamentals}

A \textbf{Discrete-Event Simulation (DES)} models the operation of a system as a discrete chronological sequence of events. Each event occurs at an instantaneous point in simulation time and initiates a state change in the network (e.g., packet arrival, backoff counter expiration, frame collision, or acknowledgment reception). Between consecutive events, no state changes occur, allowing simulation time to jump instantaneously to the next scheduled event timestamp.

\subsection{Core Components of a DES Engine}
Every discrete-event simulation engine consists of five fundamental architectural elements:
\begin{itemize}
    \item \textbf{Simulation Clock:} A variable maintaining the current virtual simulation time ($t_{\text{now}}$).
    \item \textbf{Future Event List (FEL):} A priority queue of scheduled event records ordered chronologically by execution time:
    \begin{equation}
        \text{FEL} = \left\{ (t_1, e_1), (t_2, e_2), \dots, (t_m, e_m) \right\}, \quad t_{\text{now}} \le t_1 \le t_2 \le \dots \le t_m
    \end{equation}
    Event insertion and extraction typically employ balanced binary trees (e.g., red-black trees) or calendar queues achieving $O(1)$ to $O(\log m)$ complexity.
    \item \textbf{State Variables:} Data structures encoding the network status (e.g., channel idle/busy state, buffer occupancies, contention window stages, node coordinates).
    \item \textbf{Event Routines:} Subroutines executed upon event departure that update state variables and schedule future events into the FEL.
    \item \textbf{Statistical Counters \& Trace Collection:} Data accumulators tracking performance metrics such as throughput, delay, jitter, and collision ratios.
\end{itemize}

\subsection{Statistical Rigor: Warm-up Periods and Confidence Intervals}
Because network simulations are driven by pseudo-random number generators (PRNGs), simulation outputs are random variables subject to statistical variance. Valid scientific performance evaluation requires rigorous experimental design:
\begin{enumerate}
    \item \textbf{Transient Removal (Warm-up Period):} To evaluate steady-state metrics, the simulation must run through an initial transient period $T_{\text{warm}}$ (during which empty network buffers fill to equilibrium) before accumulating statistical tallies.
    \item \textbf{Independent Replications:} The simulation is executed across $M$ independent replications using distinct PRNG seeds. For a metric of interest $X$ with sample mean $\bar{X}$ and sample variance $S^2$:
    \begin{equation}
        \bar{X} = \frac{1}{M} \sum_{k=1}^M X_k, \quad S^2 = \frac{1}{M-1} \sum_{k=1}^M (X_k - \bar{X})^2
    \end{equation}
    The two-sided $(1 - \alpha) \times 100\%$ confidence interval is given by:
    \begin{equation}
        \left[ \bar{X} - t_{M-1, 1 - \alpha/2} \frac{S}{\sqrt{M}}, \;\; \bar{X} + t_{M-1, 1 - \alpha/2} \frac{S}{\sqrt{M}} \right]
    \end{equation}
    where $t_{M-1, 1 - \alpha/2}$ denotes the critical value of Student's $t$-distribution with $M - 1$ degrees of freedom.
\end{enumerate}

Table~\ref{tab:sim_tools_overview} provides a comparative taxonomy of the major simulation environments analyzed in this appendix.

\begin{table}[htbp]
\centering
\caption{Taxonomy and comparison of network simulation and modeling platforms.}
\label{tab:sim_tools_overview}
\resizebox{\textwidth}{!}{
\begin{tabular}{lccccc}
\toprule
\textbf{Attribute} & \textbf{MATLAB} & \textbf{NS-2} & \textbf{NetSim} & \textbf{SimPy} & \textbf{NS-3} \\
\midrule
\textbf{Execution Paradigm} & Numerical / Matrix & Event-Driven & Event-Driven & Process-Driven & Event-Driven \\
\textbf{Core Engine Language} & C++ / Native & C++ & C & Python & C++ \\
\textbf{Scripting / Front-End} & MATLAB / Simulink & OTcl & GUI / XML / C & Python & C++ / Python \\
\textbf{PHY Layer Fidelity} & Waveform / Bit-level & Empirical SINR & Analytical / BER & Abstract / Error & Analytical / SINR \\
\textbf{MAC Layer Customization} & High (Scripted) & High (C++ modifications) & High (C code) & Very High (Pure Python) & High (C++ modules) \\
\textbf{Emulation / Live Traffic} & Instrument Control & Network Animator (NAM) & Hardware-in-Loop (HIL) & Real-time Environment & Raw Socket Tap \\
\textbf{Learning Curve} & Low & Steep (Dual language) & Low to Medium & Low & Moderate to Steep \\
\bottomrule
\end{tabular}
}
\end{table}

---

\section{MATLAB and Simulink for Wireless Modeling}
\label{sec:matlab_tools}

MATLAB (The MathWorks, Inc.) is the standard computational platform for mathematical modeling, matrix numerical algebra, and physical-layer waveform simulation \cite{matlab_wlan}. In the context of wireless network performance evaluation, MATLAB serves two primary roles: (i) solving numerical equations arising from analytical models (e.g., Markov chain steady-state vectors, fixed-point equations, and fluid-flow differential equations), and (ii) executing Monte Carlo simulations of medium access contention.

\subsection{Solving Non-Linear Fixed-Point Systems and Fluid Flow ODEs}
Throughout this monograph, analytical MAC models culminate in non-linear coupled systems that lack closed-form algebraic inversions. For example, in Chapter~1 and Chapter~2, Bianchi's slotted CSMA/CA throughput requires solving the coupled non-linear system:
\begin{equation}
    \tau = \frac{2(1 - 2p)}{(1 - 2p)(W + 1) + p W (1 - (2p)^m)}, \quad p = 1 - (1 - \tau)^{n-1}
\end{equation}
In MATLAB, such systems are solved to machine precision via multidimensional root-finding routines using the Levenberg-Marquardt or trust-region-dogleg algorithms:
\begin{lstlisting}[language=Matlab]
function tau = solve_bianchi(n, W, m)
    % Objective function F([tau, p]) = 0
    bianchi_eqs = @(x) [
        x(1) - (2*(1 - 2*x(2))) / ((1 - 2*x(2))*(W + 1) + x(2)*W*(1 - (2*x(2))^m));
        x(2) - (1 - (1 - x(1))^(n - 1))
    ];
    opts = optimoptions('fsolve', 'Display', 'off', 'TolFun', 1e-12);
    sol  = fsolve(bianchi_eqs, [0.05; 0.05], opts);
    tau  = sol(1);
end
\end{lstlisting}

Similarly, the non-stationary fluid flow approximations formulated in Chapter~\ref{sec:chapter_sdn} (Eqs.~\ref{eq:psffa_switch}--\ref{eq:psffa_controller}) are numerically integrated using adaptive Runge-Kutta Dormand-Prince pairs (`ode45`) or implicit stiff solvers (`ode15s`):
\begin{lstlisting}[language=Matlab]
function dydt = sdn_fluid_ode(t, y, mu_s, mu_c, q, alpha_func)
    xs = max(y(1), 0); xc = max(y(2), 0);
    alpha = alpha_func(t);
    dxs = -mu_s * (xs / (xs + 1)) + alpha + mu_c * (xc / (xc + 1));
    dxc = -mu_c * (xc / (xc + 1)) + q * alpha;
    dydt = [dxs; dxc];
end
\end{lstlisting}

\subsection{Monte Carlo MAC Simulation in MATLAB}
MATLAB supports event-by-event Monte Carlo simulation of wireless contention. In a slotted random-access environment, time is discretized into slot epochs, and node backoff counters are represented as vector operations. Listing~\ref{lst:matlab_csma} provides a complete, vector-optimized MATLAB implementation of a multi-station slotted CSMA/CA contention process tracking collisions, backoff decrements, and throughput.

\begin{lstlisting}[language=Matlab, caption={Vectorized Monte Carlo simulation of slotted CSMA/CA in MATLAB.}, label={lst:matlab_csma}]
function S = simulate_csma(N_nodes, CW_min, CW_max, max_slots)
    % State: backoff counter and contention window stage per node
    CW = repmat(CW_min, N_nodes, 1);
    BC = randi([0, CW_min - 1], N_nodes, 1);
    success_slots = 0; collision_slots = 0; idle_slots = 0;
    
    for slot = 1:max_slots
        transmitting = find(BC == 0);
        k = length(transmitting);
        if k == 0
            % Idle slot: all active nodes decrement backoff
            idle_slots = idle_slots + 1;
            BC = max(BC - 1, 0);
        elseif k == 1
            % Successful transmission
            success_slots = success_slots + 1;
            tx_node = transmitting;
            CW(tx_node) = CW_min; % Reset contention window
            BC(tx_node) = randi([0, CW_min - 1]);
            % Other nodes freeze or decrement based on protocol rules
            non_tx = setdiff(1:N_nodes, tx_node);
            BC(non_tx) = max(BC(non_tx) - 1, 0);
        else
            % Collision: k >= 2 nodes transmitted simultaneously
            collision_slots = collision_slots + 1;
            for idx = 1:k
                col_node = transmitting(idx);
                CW(col_node) = min(CW(col_node) * 2, CW_max);
                BC(col_node) = randi([0, CW(col_node) - 1]);
            end
            non_tx = setdiff(1:N_nodes, transmitting);
            BC(non_tx) = max(BC(non_tx) - 1, 0);
        end
    end
    S = success_slots / max_slots; % Normalized throughput
end
\end{lstlisting}

\subsection{Communications, WLAN, and 5G Toolboxes}
For high-fidelity physical-layer modeling, MATLAB provides specialized standard-compliant toolboxes:
\begin{itemize}
    \item \textbf{WLAN Toolbox \cite{matlab_wlan}:} Implements end-to-end IEEE 802.11a/b/g/n/ac/ax/be waveforms, multi-user OFDMA tone mapping, resource unit (RU) configurations, MIMO spatial beamforming, and TGax/TGbe fading channel models.
    \item \textbf{5G Toolbox:} Generates 3GPP Release 15/16/17 5G NR uplink and downlink waveforms, Synchronization Signal Blocks (SSB), physical random access channels (PRACH), and low-density parity-check (LDPC) / Polar channel coding.
\end{itemize}

---

\section{Network Simulator 2 (NS-2)}
\label{sec:ns2_architecture}

The Network Simulator 2 (NS-2) is historically one of the most widely cited open-source discrete-event network simulators in wireless networking research \cite{fall2003ns}. Developed through the VINT collaboration (UC Berkeley, LBL, USC/ISI, and Xerox PARC), NS-2 introduced a split-programming model designed to decouple protocol implementation speed from simulation configuration flexibility.

\subsection{The Dual-Language Architecture (C++ and OTcl)}
NS-2 separates its codebase into two distinct languages linked by an object-oriented binding mechanism called **TclCL** (Tcl with C++ Classes), as illustrated in Figure~\ref{fig:ns2_arch}:
\begin{enumerate}
    \item \textbf{Compiled C++ Hierarchy (Data Plane):} Implements packet structures, state machines, buffer queues, MAC contention timers, and PHY-layer models where high execution speed and low memory overhead are mandatory.
    \item \textbf{Interpreted Object Tcl (OTcl) Hierarchy (Control Plane):} Serves as the configuration and scripting language. Researchers write OTcl scripts to assemble topologies, instantiate node agents, bind protocols, configure traffic sources (CBR, FTP), schedule simulation events, and direct output traces without recompiling the simulator binary.
\end{enumerate}

\begin{figure}[htbp]
\centering
\begin{minipage}{0.85\textwidth}
\begin{lstlisting}[language=C++]
// Conceptual C++ / TclCL binding in NS-2
class Mac802_11Class : public TclClass {
public:
    Mac802_11Class() : TclClass("Mac/802_11") {}
    TclObject* create(int, const char*const*) {
        return (new Mac802_11());
    }
} class_mac802_11;
\end{lstlisting}
\end{minipage}
\caption{TclCL binding mechanism connecting C++ protocol objects to interpreted OTcl classes in NS-2 \cite{fall2003ns}.}
\label{fig:ns2_arch}
\end{figure}

\subsection{Wireless Networking Architecture in NS-2}
The mobile wireless extension in NS-2 (contributed by the Monarch project at Carnegie Mellon University) organizes mobile nodes into structured protocol stacks:
\begin{itemize}
    \item \textbf{MobileNode Object:} Represents a physical wireless device, encapsulating position coordinates $(X, Y, Z)$, mobility trajectories, and network interfaces.
    \item \textbf{Channel Model (`Channel/WirelessChannel'):} Propagates radio signals across all receiving interfaces within spatial range.
    \item \textbf{Radio Propagation Model:} Computes received signal power $P_r$ based on transmitter-receiver distance $d$. Implemented models include Free Space ($P_r \propto 1/d^2$), Two-Ray Ground ($P_r \propto 1/d^4$ for $d > d_c$), and Shadowing (log-normal random fading).
    \item \textbf{Physical Layer (`Phy/WirelessPhy'):} Compares received power against two physical thresholds:
    \begin{align}
        P_r &\ge \text{CSThresh\_} \implies \text{Carrier Sense Triggered (Medium Busy)} \\
        P_r &\ge \text{RXThresh\_} \implies \text{Successful Packet Reception Initiated}
    \end{align}
    \item \textbf{MAC Layer (`Mac/802\_11'):} Full state-machine implementation of IEEE 802.11 DCF, complete with DIFS/SIFS deferral, physical/virtual carrier sensing (NAV), slotted binary exponential backoff, RTS/CTS handshakes, and retry counters.
\end{itemize}

\subsection{Simulation Scripting in OTcl}
Listing~\ref{lst:ns2_script} demonstrates a canonical NS-2 wireless simulation script setting up an IEEE 802.11 ad-hoc network with Two-Ray Ground propagation and UDP/CBR traffic.

\begin{lstlisting}[language=tcl, caption={NS-2 OTcl configuration script for a mobile wireless ad-hoc network.}, label={lst:ns2_script}]
# Instantiate Simulator and Event Scheduler
set ns_ [new Simulator]
set tracefd [open wireless_out.tr w]
$ns_ trace-all $tracefd

# Configure Mobile Node Protocol Stack
set topo [new Topography]
$topo load_flatgrid 500 500
create-god 2

$ns_ node-config -adhocRouting DSDV \
                 -llType LL \
                 -macType Mac/802_11 \
                 -ifqType Queue/DropTail/PriQueue \
                 -ifqLen 50 \
                 -antType Antenna/OmniAntenna \
                 -propType Propagation/TwoRayGround \
                 -phyType Phy/WirelessPhy \
                 -channelType Channel/WirelessChannel \
                 -topoInstance $topo

# Instantiate Wireless Nodes
set node_(0) [$ns_ node]
set node_(1) [$ns_ node]
$node_(0) set X_ 50.0; $node_(0) set Y_ 100.0; $node_(0) set Z_ 0.0
$node_(1) set X_ 200.0; $node_(1) set Y_ 100.0; $node_(1) set Z_ 0.0

# Attach CBR Traffic over UDP
set udp [new Agent/UDP]
$ns_ attach-agent $node_(0) $udp
set null [new Agent/Null]
$ns_ attach-agent $node_(1) $null
$ns_ connect $udp $null

set cbr [new Application/Traffic/CBR]
$cbr set packetSize_ 1000
$cbr set rate_ 1Mb
$cbr attach-agent $udp

# Schedule Events
$ns_ at 1.0 "$cbr start"
$ns_ at 10.0 "$cbr stop"
$ns_ at 10.1 "finish"
$ns_ run
\end{lstlisting}

\subsection{Trace Analysis and Post-Processing}
NS-2 logs all packet transitions into an ASCII trace file (`.tr`). In wireless trace format, each entry records the event type (`s` for send, `r` for receive, `d` for drop, `f` for forward), timestamp, node ID, trace layer (`MAC`, `RTR`, `AGT`), packet sequence number, packet size, and MAC-layer duration flags. Extraction of performance metrics is executed via AWK scripts or Python data parsers:
\begin{lstlisting}[language=bash]
# Compute End-to-End Packet Delivery Ratio (PDR) via AWK
awk '
BEGIN { sent = 0; received = 0; }
$1 == "s" && $4 == "AGT" { sent++; }
$1 == "r" && $4 == "AGT" { received++; }
END { printf("PDR = %.2f%%\n", (received / sent) * 100); }
' wireless_out.tr
\end{lstlisting}

---

\section{NetSim (by Tetcos)}
\label{sec:netsim_platform}

NetSim, developed by Tetcos (Bangalore, India), is a leading commercial discrete-event network simulator and emulator widely utilized in academia, defense research, and telecommunications R\&D \cite{tetcos2020netsim}. Unlike open-source research platforms that often require low-level patch maintenance, NetSim provides an integrated GUI environment, validated protocol libraries, packet-level animation, and real-time Hardware-in-the-Loop (HIL) emulation capabilities.

\subsection{Architecture and Component Hierarchy}
NetSim is architected around a central discrete-event simulation kernel implemented in C. Its modular engine is organized hierarchically:
\begin{itemize}
    \item \textbf{Simulation Kernel (Engine):} Maintains the global simulation timeline, calendar event list, and memory management pools.
    \item \textbf{Protocol Libraries:} Pre-built, standard-compliant implementations of network protocols compiled as modular dynamic link libraries (DLLs). Supported technologies include:
    \begin{itemize}
        \item \textit{WLAN:} IEEE 802.11a/b/g/n/ac/ax/e (EDCA, MU-MIMO, OFDMA).
        \item \textit{IoT \& WSN:} IEEE 802.15.4 (CSMA/CA and TSCH), 6LoWPAN, RPL, and CoAP.
        \item \textit{Cellular Networks:} 4G LTE, LTE-Advanced, and full 5G NR (mmWave beamforming, numerologies, gNB/UE dual connectivity, 5G Core).
        \item \textit{Cognitive Radio:} IEEE 802.22 WRAN, spectrum sensing, and primary user protection.
    \end{itemize}
    \item \textbf{Network Design Interface:} Graphical CAD canvas enabling drag-and-drop network topology construction, link parameterization, and node configuration.
    \item \textbf{Metrics Engine:} Automatically aggregates link-level and application-level statistics (throughput, delay, jitter, queue backlogs, bit error rates) into structured HTML/XML reports.
\end{itemize}

\subsection{The Four-Event Discrete-Event Execution Model}
At the core of NetSim's MAC and physical-layer engines is a standardized four-event lifecycle governing packet transmission:
\begin{enumerate}
    \item \textbf{PHYSICAL\_IN\_EVENT:} Triggered when a packet signal begins arriving at a physical interface from the wireless channel. The engine initiates interference calculations, accumulating overlapping signals into the instantaneous SINR calculation.
    \item \textbf{PHYSICAL\_OUT\_EVENT:} Triggered when packet transmission across the radio interface completes, releasing transceiver hardware resources and notifying the channel module.
    \item \textbf{MAC\_IN\_EVENT:} Executed when an upper-layer network packet arrives at the MAC transmission buffer. The MAC engine checks buffer occupancy, verifies contention state, and schedules backoff timers.
    \item \textbf{MAC\_OUT\_EVENT:} Executed when medium contention succeeds (e.g., backoff counter reaches zero and channel is sensed idle for DIFS/AIFS). The MAC prepends headers, calculates frame duration, and pushes the MPDU to the physical layer.
    \item \textbf{TIMER\_EVENT:} Manages protocol-specific timers including SIFS timers, ACK timeout timers, beacon broadcast periods, and TSCH time-slot boundaries.
\end{enumerate}

\subsection{Advanced Emulation and Hardware-in-the-Loop (HIL)}
A distinguished feature of NetSim is its **Emulation Server**. NetSim can bind simulated virtual nodes to physical Network Interface Cards (NICs) on the host workstation. Real physical IP traffic (generated by physical laptops, IoT sensors, or video streaming clients) enters the NetSim kernel, traverses a simulated multi-hop wireless mesh with configurable fading, path loss, and congestion delay, and exits back to physical receiver endpoints. This enables researchers to evaluate how real-world software applications respond to simulated next-generation MAC protocols.

Furthermore, NetSim provides native co-simulation interfaces with **MATLAB** and **Wireshark**. Packet headers captured during simulation runs can be streamed dynamically to Wireshark for deep frame-level inspection, while complex beamforming and channel estimation routines can be delegated to MATLAB via native C-MEX bindings.

---

\section{SimPy: Process-Oriented Simulation in Python}
\label{sec:simpy_framework}

SimPy (Simulation in Python) is an open-source, process-oriented discrete-event simulation framework implemented in native Python \cite{matloff2008introduction}. Unlike thread-based simulation environments that incur substantial operating system context-switching overhead, SimPy leverages **Python generator functions** (`yield`) to model concurrent active components as lightweight, cooperative coroutines.

\subsection{The Process-Oriented Paradigm and Core Primitives}
In SimPy, simulation entities (e.g., mobile stations, wireless channels, access points) are modeled as continuous Python processes. The core primitives include:
\begin{itemize}
    \item \textbf{\texttt{simpy.Environment}}: Manages simulation time and the event queue. The environment advances time when a process yields an event.
    \item \textbf{\texttt{env.timeout(duration)}}: Suspends the calling process for a specified virtual duration. This models physical operations such as frame transmission airtime, propagation delay, or backoff slot countdowns.
    \item \textbf{\texttt{simpy.Resource}}: Models shared capacity-constrained resources (e.g., shared wireless channels, processing servers). Processes request access and release resources upon completion.
    \item \textbf{\texttt{simpy.Container}} / \textbf{\texttt{simpy.Store}}: Models finite packet buffers, producer-consumer queues, and message mailboxes.
\end{itemize}

\subsection{MAC Protocol Evaluation using SimPy: Contributions of the IITM Mobile Computing Lab}
In 2010, C.~Hemanth, T.~G.~Venkatesh, and Sriram V.~Iyer pioneered the application of SimPy for evaluating next-generation wireless medium access control protocols \cite{hemanth2010simpy}. They demonstrated that while C++ simulators (such as NS-2) offered raw computational speed, their complex class hierarchies and steep development curves significantly slowed protocol experimentation. By modeling wireless stations as concurrent SimPy processes, they achieved dramatic reductions in code complexity while maintaining rigorous fidelity to 802.11 backoff dynamics and directional 60~GHz millimeter-wave scheduling \cite{hemanth2010simpy}.

Listing~\ref{lst:simpy_mac} illustrates an idiomatic, modern SimPy implementation of a multi-station CSMA/CA MAC protocol featuring carrier sensing, random backoff, collision detection, and throughput measurement.

\begin{lstlisting}[language=Python, caption={SimPy implementation of CSMA/CA contention with backoff and collision detection \cite{hemanth2010simpy}.}, label={lst:simpy_mac}]
import simpy
import random

class WirelessChannel:
    """Models a shared wireless broadcast channel."""
    def __init__(self, env):
        self.env = env
        self.active_transmissions = 0
        self.collision = False

    def is_idle(self):
        return self.active_transmissions == 0

class WirelessStation:
    """Models a contending CSMA/CA node with exponential backoff."""
    def __init__(self, env, node_id, channel, cw_min=16, cw_max=1024, pkt_time=1.0, slot_time=0.02):
        self.env = env
        self.id = node_id
        self.channel = channel
        self.cw_min = cw_min
        self.cw_max = cw_max
        self.cw = cw_min
        self.pkt_time = pkt_time
        self.slot_time = slot_time
        self.successful_pkts = 0
        self.action = env.process(self.run())

    def run(self):
        while True:
            # 1. Generate backoff counter uniformly in [0, CW - 1]
            backoff_slots = random.randint(0, self.cw - 1)
            
            # 2. Countdown backoff counter when channel is idle
            while backoff_slots > 0:
                yield self.env.timeout(self.slot_time)
                if self.channel.is_idle():
                    backoff_slots -= 1
            
            # 3. Initiate Transmission
            self.channel.active_transmissions += 1
            if self.channel.active_transmissions > 1:
                self.channel.collision = True
            
            # Hold channel for packet transmission duration
            yield self.env.timeout(self.pkt_time)
            
            # 4. Check Transmission Outcome
            if not self.channel.collision:
                # Successful transmission
                self.successful_pkts += 1
                self.cw = self.cw_min  # Reset contention window
            else:
                # Collision occurred: double contention window
                self.cw = min(self.cw * 2, self.cw_max)
            
            self.channel.active_transmissions -= 1
            if self.channel.active_transmissions == 0:
                self.channel.collision = False  # Reset channel collision state

# Simulation Setup and Execution
def run_simulation(num_nodes=10, sim_time=1000.0):
    env = simpy.Environment()
    channel = WirelessChannel(env)
    stations = [WirelessStation(env, i, channel) for i in range(num_nodes)]
    env.run(until=sim_time)
    
    total_delivered = sum(s.successful_pkts for s in stations)
    throughput = (total_delivered * 1.0) / sim_time
    print(f"Nodes: {num_nodes} | Delivered: {total_delivered} | Channel Throughput: {throughput:.3f}")
    return throughput
\end{lstlisting}

\subsection{Advantages and Trade-Offs of SimPy}
The primary advantages of SimPy in wireless networking research include:
\begin{itemize}
    \item \textbf{Expressive Elegance:} Protocols are written in linear, procedural Python syntax using native language constructs, eliminating nested state-machine callbacks.
    \item \textbf{Scientific Ecosystem Integration:} Seamlessly couples with NumPy for numerical linear algebra, SciPy for optimization, Matplotlib/Seaborn for publication-grade visualization, and PyTorch/TensorFlow for reinforcement-learning-driven MAC scheduling (such as the deep RL algorithms in Chapter~\ref{sec:chapter_ml}).
    \item \textbf{Rapid Prototyping:} Ideal for conceptual validation of novel MAC mechanisms (e.g., multi-packet reception, SIC ordering, dynamic slicing) before committing to complex C++ simulator implementations.
\end{itemize}
The principal trade-off is execution speed: pure Python interpretation executes slower than compiled C/C++ engines for large-scale networks with tens of thousands of simultaneous nodes. However, for topologies with tens to hundreds of nodes, SimPy delivers outstanding development speed and analytical transparency.

---

\section{Network Simulator 3 (NS-3) and Comparative Assessment}
\label{sec:ns3_summary}

Network Simulator 3 (NS-3) is the modern successor to NS-2, engineered from the ground up as a pure C++ library with optional Python bindings \cite{henderson2008network}. Unlike NS-2, NS-3 completely abandons OTcl, eliminating the architectural impedance mismatch of the dual-language model.

\subsection{Key Architectural Innovations of NS-3}
\begin{itemize}
    \item \textbf{Exact Byte-Level Packet Buffers:} In NS-3, the `ns3::Packet` class manages an exact serialized byte buffer alongside a zero-copy metadata tag system (`ns3::Header`, `ns3::Trailer`). This allows packets to be serialized into real PCAP files that can be examined with Wireshark.
    \item \textbf{Attribute and Configuration System:} Every protocol parameter in NS-3 is registered through an object attribute system, enabling command-line parameter overrides and XML export without code modification.
    \item \textbf{Direct Code Execution (DCE):} Enables unmodified POSIX C/C++ application binaries (e.g., Linux kernel network stacks, routing daemons) to run directly within the simulated virtual environment.
    \item \textbf{State-of-the-Art Wi-Fi \& 5G Modules:} The \texttt{wifi} module in NS-3 provides the most rigorously validated implementation of IEEE 802.11ax/be available in open-source software, including dynamic OFDMA sub-channel scheduling, Multi-User MIMO, and spatial reuse (BSS color filtering). The \texttt{nr} (5G-LENA) module supports full 3GPP NR Release 16 features.
\end{itemize}

\subsection{Comparative Summary and Tool Selection Matrix}
Selecting the appropriate tool depends directly on the research objective, as summarized in Table~\ref{tab:tool_selection_guide}.

\begin{table}[htbp]
\centering
\caption{Practical selection guide for wireless performance evaluation tools.}
\label{tab:tool_selection_guide}
\resizebox{\textwidth}{!}{
\begin{tabular}{lp{6cm}p{7cm}}
\toprule
\textbf{Tool} & \textbf{Recommended Research Applications} & \textbf{Key Advantages} \\
\midrule
\textbf{MATLAB} & 
\begin{itemize}
    \item Solving non-linear analytical equations (Bianchi, PSFFA).
    \item Bit-level PHY waveform and link-level BER simulations.
    \item Rapid Monte Carlo prototyping of new backoff algorithms.
\end{itemize} & 
Comprehensive mathematical libraries, instant matrix plotting, validated standard toolboxes (WLAN, 5G). \\
\midrule
\textbf{NS-2} & 
\begin{itemize}
    \item Replicating legacy wireless ad-hoc and sensor network literature.
    \item Routing and MAC protocol cross-comparison.
\end{itemize} & 
Massive historical literature base, extensive community patches for ad-hoc protocols. \\
\midrule
\textbf{NetSim} & 
\begin{itemize}
    \item Enterprise network planning and 5G/IoT campus deployment.
    \item Hardware-in-the-Loop (HIL) physical device emulation.
    \item Wireshark-level packet protocol verification.
\end{itemize} & 
Commercial GUI, out-of-the-box standard compliance, verified protocol interaction, direct tech support. \\
\midrule
\textbf{SimPy} & 
\begin{itemize}
    \item Algorithmic MAC protocol prototyping (MPR, SIC, TSCH).
    \item Integrating machine learning / RL with MAC scheduling.
    \item Transient queueing and game-theoretic network simulations.
\end{itemize} & 
Clean process-oriented Python syntax, fast development cycles, seamless integration with AI/ML libraries. \\
\midrule
\textbf{NS-3} & 
\begin{itemize}
    \item High-fidelity large-scale 5G NR and Wi-Fi 6/7 MAC/PHY evaluations.
    \item Full-stack end-to-end transport (TCP/QUIC) performance over wireless.
\end{itemize} & 
Pure C++ speed, byte-level packet serialization, Linux DCE integration, active standard evolution. \\
\bottomrule
\end{tabular}
}
\end{table}

The simulation environments detailed in this appendix provide the empirical counterpart to the analytical formulations developed across the monograph, allowing researchers to evaluate theoretical boundaries under realistic operational constraints.
